\vspace{-0.3cm}
\section{Introduction}
\label{sec:intro}


Large language models~(LLMs) have exhibited remarkable performance across various natural language processing tasks~\cite{openai2023chatgpt,yang2024qwen2,dubey2024llama,anthropic2024introducing}.
Concurrently, transformers for visual tasks~(\eg, ViT~\cite{dosovitskiy2021an}) and cross-modal visual-language models~(VLMs), 
such as CLIP~\cite{radford2021learning} and BLIP~\cite{li2022blip}, have advanced rapidly.
Riding the waving of LLMs, multimodal large language models~(MLLMs) are transforming the landscape of computer vision globally~\cite{wang2024videoagent,Qwen2VL,lu2024ovis,wang2024videotree,TanKoala2024}. 


MLLMs typically use the visual encoder from VLMs and a new-trained projector to extract a long sequence of visual tokens from images.
The combined sequence of visual and text tokens is then fed into LLMs for visual understanding~\cite{liu2023llava,cheng2024videollama2,li2023mvbench}.
However, a video contains hundreds or even thousands of frames. 
Directly concatenating sequences from multiple frames and feeding them into Video-LLMs poses several challenges: 
1)~LLMs' limited context length can not accommodate excessively long visual sequences;
2)~Redundant and duplicate frames can hinder visual perception, as irrelevant frames may cause ``lost-in-the-middle'' issues~\cite{liu2024st,yu2024frame}.
3)~Edge-deployed LLMs face resource constraints when processing overly long input sequence, and
4)~Proprietary models charge based on token usage, making excessive frame input costly.

\begin{figure}[t]
	\centering
  	\includegraphics[width=0.8\linewidth]{figures/llava_ov_video_mme.pdf}
	\caption{
	    Overall accuracy (\%) on Video-MME~(without subtitles) for LLaVA-OneVision, 
    	comparing uniform sampling and \mname{}\ selected from 128 frame candidates.
% 		Accuracy (\%) on Video-MME~(w/o subs) for LLaVA-OneVision, uniform sampling and \mname{}\ selected from 128frm.
	}
	\label{fig:gain_on_videomme}
\end{figure}


As shown in the empirical example in \cref{fig:gain_on_videomme}, 
the blue line illustrates that the marginal benefits of increasing the number of frames diminish, 
with performance even declining when using 64 frames instead of 32 frames. 
This indicates a law of diminishing marginal utility as the number of frames increases, 
and additional frames may even have negative effects. 
Hence, selecting a limited number of key frames is crucial for effective and efficient video understanding. 



In previous research, uniform frame sampling is a common strategy in Video-LLMs~\cite{miao2023towards,lin2023video,zhang2024video,wang2024internvideo2}. 
However, it is arbitrary, especially in video question-answering~(VidQA) tasks, 
as it disregards the current query, potentially missing important key frames and including irrelevant ones. 
Furthermore, \citet{liang2024keyvideollm,wang2024weakly,yu2024self} have explored using the query to retrieve the most relevant frames.
However, these point-wise frame selection methods overlook the list-wise relationships among selected frames, including: 
1)~\emph{Diversity}: Videos often contain many similar frames; redundantly selecting them increases the computational burden during inference with little information gain and may hinder the usage of text information.
2)~\emph{Sequentiality}: The inherent sequential nature of frames is essential. 
For instance, when asking, ``What did the chef take \underline{after} adding the spices?'' in a cooking video, frames depicting the steps following the addition of the spices are essential.



Building on the insights from the above analysis, we emphasize that effective frame selection in Video-LLMs should adhere to three key principles:
\textbf{query relevance, list-wise diversity, and sequentiality.} 
Prior research has addressed some of them but has not fully encompassed all of them, particularly the latter two~\cite{jin2024chat,wang2025ViLA,tang2025,hu2025m}.
Modeling list-wise relationships in subset frame selection presents two main challenges:
1)~the lack of available supervised list-wise datasets; 
2)~even if a dataset is constructed to learn a list-wise scoring function, finding the optimal subset with the highest score remains an NP-hard problem that cannot be solved in polynomial time complexity~\cite{bian2020efficient,chen2018fast}.



In this paper, we propose a training-free, model-agnostic list-wise frame selection method called 
\underline{M}arkov~\underline{D}ecision \underline{D}eterminantal~\underline{P}oint \underline{P}rocess~with \underline{D}ynamic~\underline{P}rogramming, abbreviated as \mname{}.
We reuse pretrained VLMs to compute list-wise score for the subset using Determinantal Point Process~(DPP), 
thereby eliminating the need for collecting a supervised dataset or training a scoring model from scratch.
DPP was originally introduced to characterize the Pauli exclusion principle, where no two fermions can occupy the same quantum state.
This repulsion, interpreted as list-wise diversity, is effectively modeled by DPP~\cite{macchi1975coincidence},
which can obtain a $(1-1/e)$-approximate solution with polynomial complexity~\cite{chen2018fast}.
%
In \mname{}, we extend standard DPP to consider query relevance, list-wise diversity, and sequentiality for frame selection in VidQA.



Specifically, we define a conditional multiple Gaussian kernel~(CMGK) to estimate high-dimensional frame similarities conditioned on the query within the reproducing kernel Hilbert space~(RKHS), rather than using a query-agnostic low-dimensional measure such as cosine similarity. 
We then apply DPP to the conditional similarity matrix generated by CMGK to capture the relevance of frames to the current query and ensures list-wise diversity among them. 
To incorporate sequentiality, we partition the video into continuous segments and apply DPP within each segment, conditioned on the selection from the immediately preceding one.
This process is modeled as a Markov Decision Process~(MDP) to optimize the allocation of the total selection size $k$ across segments, resulting in a solution with pseudo-polynomial time complexity.


To evaluate \mname{}, we integrate it as a plug-and-play process at the test time of state-of-the-art~(SOTA) Video-LLMs, including VILA-V1.5~\cite{lin2024vila}, MiniCPM-V2.6~\cite{yao2024minicpm}, LLaVA-OneVision~\cite{li2024llava}, and Ovis2~\cite{ovis2}.
We conduct extensive experiments on three widely used long-video benchmarks: Video-MME~\cite{fu2024video}, MLVU~\cite{MLVU}, and LongVideoBench~\cite{wu2024longvideobench}.
Experimental results show significant performance improvements over baseline methods. 


Furthermore, we summarize our main contributions:
\begin{enumerate}
    \item We demonstrate that effective frame selection should account for query relevance, list-wise diversity, and sequentiality. The proposed \mname{} is the first to address all three aspects comprehensively.
    \item \mname{} is effective, offering a \((1-1/e)\)-approximate solution to the NP-hard list-wise subset frame selection problem, achieving significant performance improvements over existing baselines.
    \item \mname{} is efficient, extending DPP methods by optimally allocating the selection capacity \(k\) across the video for VidQA, maintaining pseudo-polynomial time complexity, even though this allocation problem is NP-hard.
    \item \mname{} is training-free and model-agnostic, enabling seamless plug-and-play integration with various existing Video-LLMs at test time.
\end{enumerate}





